\begin{claim}{C.68}{Proved}{Relative first-kind disk centralizers and spherical frontier quotients}
Let $O\subset S^2$ be a simply connected domain whose
frontier is a nondegenerate locally connected continuum.
Retain a Riemann coordinate and its actual continuous
Carath\'eodory quotient
\[
 \phi:\mathbb D\longrightarrow O,\qquad
 q:\overline{\mathbb D}\longrightarrow\overline O .
                                                               \tag{C.68.1}
\]
Writing $\mathcal D_\partial$ as in C.67, put
\[
 \mathcal K_O^+
  =\{b\in\operatorname{Homeo}^+(S^2):
                         b|_{S^2\setminus O}=\id\}.
                                                               \tag{C.68.2}
\]
The map $T_q$, equal to $\phi a\phi^{-1}$ on $O$
and to the identity off $O$, is an isomorphism of
topological groups
\[
                  T_q:\mathcal D_\partial
                              \xrightarrow{\ \cong\ }\mathcal K_O^+.
                                                               \tag{C.68.3}
\]
The full group $\mathcal K_O^+$ is contractible.
No injectivity of $q|_{S^1}$ is assumed.

There is also the following operator comparison.
Let $A\leq\operatorname{Homeo}(\overline{\mathbb D})$,
and retain its actual restriction closure
\[
 Q=\overline{\{a|_{\mathbb D}:a\in A\}}
             \ \leq\operatorname{Homeo}(\mathbb D).
                                                               \tag{C.68.4}
\]
Suppose that $Q$ is locally compact and acts freely
and properly on $\mathbb D$. Retain a specified
$k\in\operatorname{Homeo}(\overline{\mathbb D})$ and
a torsion-free Fuchsian group $\Lambda$ of the first
kind, with the actual subgroup inclusion
\[
       (k|_{\mathbb D})^{-1}\Lambda(k|_{\mathbb D})\leq Q .
                                                               \tag{C.68.5}
\]
Here first kind means that the limit set of the
specified disk action is all of $S^1$.
Restriction is an isomorphism of topological groups
\[
 \operatorname{res}:
       C_{\mathcal D_\partial}(A)
       \xrightarrow{\ \cong\ }
       C_{\operatorname{Homeo}(\mathbb D)}(Q).
                                                               \tag{C.68.6}
\]
Every based loop in the left-hand group has a based
filling in that same group.
Neither finite generation nor cocompactness of
$\Lambda$ is required.

For original ambient operators stabilizing $O$,
take $A$ to be their actual prime-end image under
(C.68.1). Equations (C.68.3) and (C.68.6) then
give a based filling of every loop in the corresponding
relative ambient centralizer whenever the displayed
data (C.68.4)--(C.68.5) hold. This retains each original
operator and its restriction kernel.
For the plane application, retain the specified
point $\infty\notin O$; the entire filling fixes
that point and restricts to a filling through
plane homeomorphisms.
\end{claim}

\begin{proof}
\emph{The actual spherical boundary quotient.}
The locally connected Carath\'eodory theorem and
the prime-end compactification in
\cite[Sections~1 and~3.1]{C_BinderRojasYampolsky2014}
give (C.68.1). They apply in a spherical coordinate:
the conformal map may be meromorphic in a chosen
plane chart, but is a homeomorphism from the open
disk onto $O$. The boundary map
$q_\partial=q|_{S^1}$ is onto $\partial O$.

This boundary map is constant on no nonempty open
arc. Indeed, a constant value $p$ on such an arc
admits a target complex coordinate finite near $p$.
Continuity of $q$ gives a neighborhood of a smaller
closed subarc on which that coordinate of $\phi$
is holomorphic and continuous up to the arc.
Subtract its constant boundary value and straighten
the source circle locally. Schwarz reflection across
the resulting real segment gives a holomorphic
extension with a segment of zeros. The identity
theorem makes $\phi$ constant on an open subset,
and then on the connected disk, a contradiction.

We now prove the needed precise property of its
possibly nontrivial fibers. There cannot be four
cyclically ordered points $a,b,c,d$ in $S^1$ with
\[
            q(a)=q(c)=p,\qquad q(b)=q(d)=r,\qquad p\ne r .
                                                               \tag{C.68.7}
\]
Choose an embedded closed crosscut $\alpha$ joining
$a$ to $c$ and otherwise contained in $\mathbb D$.
The map $q$ embeds the interior of this crosscut,
and its interior image lies in $O$. Its only
coincident endpoint images are the point $p$.
Consequently
\[
           J=q(\alpha)\text{ is a Jordan curve in }S^2,
           \qquad J\cap\partial O=\{p\}.
                                                               \tag{C.68.8}
\]
The two components of $\mathbb D\setminus\alpha$
have connected images disjoint from $J$.
Near an interior point of $\alpha$, the map $\phi$
is a local homeomorphism and sends its two local
sides to the two local sides of $J$.
Thus the two entire component images lie in the
different complementary Jordan domains.
The cyclic order in (C.68.7) places $b$ and $d$
on the respective boundary arcs of the two source
components. Approaching them from those components
shows that $r$ lies in the closures of both
Jordan domains, and hence in $J$.
Since $r\in\partial O$, (C.68.8) gives $r=p$,
contrary to (C.68.7). This proves that distinct
fibers of $q_\partial$ do not have alternating
pairs of points.

Let $u\in\operatorname{Homeo}^+(S^1)$ satisfy
\[
                            q_\partial u=q_\partial .
                                                               \tag{C.68.9}
\]
It follows that $u=\id$. To see this, choose a
moved point $\zeta$ and an increasing real lift
in which $\zeta<u\zeta<\zeta+1$.
For every $\eta>\zeta$ sufficiently close to
$\zeta$, continuity and strict increase give
\[
                       \zeta<\eta<u\zeta<u\eta<\zeta+1 .
                                                               \tag{C.68.10}
\]
The absence of a constant boundary arc permits
such an $\eta$ with $q_\partial\eta\ne q_\partial\zeta$.
Equation (C.68.9) and (C.68.10) then give the
forbidden alternating pairs in two distinct
fibers. This proves the assertion.

\emph{The relative spherical group, its topology, and its operators.}
Every original sphere homeomorphism $h$ preserving
$O$ acts on its prime ends. This is precisely the
null-crosscut-chain functor proved in C.67:
uniform continuity on the compact sphere preserves
null chains, their inclusion relations, and their
equivalence, and the inverse gives the inverse
prime-end map. In the retained coordinate write
this map as $\widehat h$. It satisfies
\[
 \widehat h|_{\mathbb D}=\phi^{-1}h\phi,\qquad
                             q\widehat h=hq .
                                                               \tag{C.68.11}
\]
For $h\in\mathcal K_O^+$, the boundary restriction
of $\widehat h$ preserves orientation and satisfies
(C.68.9); it is therefore the identity.

For $a\in\mathcal D_\partial$, the continuous
map $qa$ is constant on every fiber of $q$:
interior fibers are singletons, and all points
in each boundary fiber are fixed by $a$.
The compact quotient $q$ therefore descends
$qa$ to a continuous self-map of $\overline O$.
The same construction for $a^{-1}$ is its inverse.
It is the identity on $\partial O$, so pasting
with identity on $S^2\setminus O$ gives $T_q(a)$.
It preserves orientation because its restriction
to $O$ does. Equations on $O$ prove injectivity
and the homomorphism identity. Equation (C.68.11)
and boundary rigidity give its inverse, proving
the algebraic bijection (C.68.3).

Uniform continuity of $q$ gives forward continuity
of $T_q$: uniform changes of $a$ and $a^{-1}$
on the closed disk give uniform changes of
$T_q(a)$ and its inverse on the sphere.
Both groups in (C.68.3) are closed Polish subgroups
of their homeomorphism groups. The left action
of $\mathcal D_\partial$ on $\mathcal K_O^+$
through $T_q$ is continuous and transitive.
The verified Polish open orbit-map theorem used
in C.26 and C.30 makes this bijective homomorphism
open. Its inverse is continuous, as claimed.
Conjugating the explicit Alexander contraction
of $\mathcal D_\partial$ in C.67 by this
isomorphism proves contractibility of the full
relative group.

For any original operator $h$ preserving $O$,
the exact interior conjugation and the exterior
identity give
\[
          T_q(\widehat h\,a\,\widehat h^{-1})
                        =h\,T_q(a)\,h^{-1}.
                                                               \tag{C.68.12}
\]
Orientation reversal by $h$ is allowed.
The kernel of $h\mapsto\widehat h$ is exactly
the kernel of restriction to $O$.
Thus (C.68.3) identifies the actual centralizers,
with the actual induced topologies and operator
labels. No equivariance of the Alexander formula
under these operators is used.

\emph{First-kind extension with the specified disk conjugacy.}
We use the original first-kind extension theorem
in \cite[Section~3.B, Proposition~3.5 and
Corollary~3.5.1, pp.~31--32]{C_Tukia1972}.
It states that a homeomorphism of the open disk
inducing a geometric isomorphism between
first-kind groups extends homeomorphically to
the closed disk. Its boundary map is determined
by the images of the attracting fixed points
of infinite-order elements. The proof constructs
the boundary map from nested intervals with
hyperbolic-axis endpoints and proves continuity
of the actual interior homeomorphism at their
limiting boundary points. The hypotheses in
that section impose no finite generation or
cocompactness requirement.

In particular, an interior homeomorphism $g$
commuting with $\Lambda$ induces its identity
group isomorphism. The boundary map supplied
by that theorem fixes every one of the stated
attracting points; they are dense for a
first-kind group. Its boundary extension is
therefore the identity:
\[
      g\in C_{\operatorname{Homeo}(\mathbb D)}(\Lambda)
        \quad\Longrightarrow\quad
      \overline g\in\mathcal D_\partial .
                                                               \tag{C.68.13}
\]
This is an extension of the given homeomorphism,
not a replacement by another representative.
Uniqueness follows from density of the open disk.

Now let $g\in C_{\operatorname{Homeo}(\mathbb D)}(Q)$.
By the actual inclusion (C.68.5), the map
$k g k^{-1}$ on the open disk commutes with
$\Lambda$. Equation (C.68.13) extends it to
the closed disk, fixing the boundary.
Conjugating this extension by the specified
closed-disk homeomorphism $k^{-1}$ gives the
unique extension $\overline g$ of $g$, still
fixing the boundary. For each $a\in A$,
the two continuous maps $\overline g\,a$
and $a\,\overline g$ agree on the dense open
disk, and hence everywhere.
Thus $\overline g\in C_{\mathcal D_\partial}(A)$.

Conversely, a member of
$C_{\mathcal D_\partial}(A)$ restricts to a
homeomorphism commuting with all restrictions
of $A$, and then with their actual closure
$Q$, by continuity of group operations.
This proves surjectivity and injectivity
of (C.68.6). Restriction is continuous.
The source and target are closed Polish
subgroups: commuting with each specified
operator is a closed condition, and fixing
the boundary is closed.
Applying the same open orbit-map theorem
to the transitive action through restriction
proves continuity of its inverse.
Consequently every compact parameter family,
including its inverse family, extends
continuously through this exact inverse map.
No separate assumption about a boundary
restriction fibration is needed for this
particular isomorphism.

\emph{Based fillings and the original frontier.}
Let $t\mapsto f_t$ be a based loop in
$C_{\mathcal D_\partial}(A)$.
As a loop in the full group
$\mathcal D_\partial$, it is null-homotopic
by the Alexander contraction in C.67.
Restriction makes it an ambient-null loop
in $\operatorname{Homeo}(\mathbb D)$.
The proper-action injection C.33 applies
to the exact group $Q$ in (C.68.4).
It supplies a based filling
\[
       B:D^2\longrightarrow
                    C_{\operatorname{Homeo}(\mathbb D)}(Q)
\]
of that restricted loop. Composing with
the continuous inverse in (C.68.6) gives
a based filling of the original $f_t$
through boundary-fixed closed-disk
homeomorphisms commuting with every
original $a\in A$.
This proves the asserted loop statement.

For the spherical domain, use
$A=\{\widehat h:h\in H_O\}$ from (C.68.11).
Apply (C.68.6) to the loop obtained through
$T_q^{-1}$, and return the resulting filling
through $T_q$. Equation (C.68.12) proves
commutation with every original operator
in $H_O$ throughout the filling; the
restriction kernel acts identically on
$O$ throughout. The complement of $O$
is fixed pointwise.
This passage uses the quotient $q$ and
never uses an inverse boundary value of
that quotient. It therefore includes slit
frontiers and other locally connected
frontiers with multiple prime-end accesses.

In the plane application $\infty\notin O$,
so $\infty$ belongs to the fixed complement.
Restriction from the sphere group fixing
$\infty$ gives the original compact-open
and inverse topology on plane homeomorphisms;
the converse continuity is the one-end
compact-core comparison proved in C.66.
Hence this is a filling in the required
plane group with its specified infinity.
\end{proof}

\emph{Scope of the extension.}
The closed-disk homeomorphism $k$ and the
first-kind subgroup inclusion (C.68.5) are
part of the proved comparison.
An arbitrary interior uniformization of
$O/Q$ does not supply this $k$ on the
original prime-end circle.
The argument makes no such substitution.
It also supplies no uniform estimate for
pasting contractions over infinitely many
different transported domains.
The remaining general frontier and
moving-quotient extension in C.12 stays open.
